\documentclass[12pt,a4paper]{article}
\usepackage[spanish,es-tabla]{babel}
\decimalpoint
\usepackage{array}
\newcolumntype{C}[1]{>{\centering\arraybackslash}m{#1}}
\usepackage{bbm}
\usepackage{bm}
\usepackage{tikz}
\usetikzlibrary{fit,positioning}
\usepackage{pgfplots}
\pgfplotsset{compat=1.18}
\usepackage{braket}
\usepackage[utf8]{inputenc}
\usepackage{graphicx}
\usepackage{enumitem}
\usepackage{subcaption}
\usepackage{cancel}
\usepackage{multicol}
\usepackage{wrapfig}
\usepackage[T1]{fontenc}
\usepackage{url}
\usepackage{graphicx}
\usepackage{gensymb}
\usepackage{booktabs}
\usepackage{amssymb, amsmath}
\usepackage{wrapfig}
\parskip= 0.5 mm
\oddsidemargin = -1 cm
\headsep = -2 cm
\textwidth = 17.5cm
\textheight = 25.5cm
\begin{document}
\begin{center}
    {\LARGE \textbf{Introducción a Física Cuántica}} \\[0.4cm]

    {\large Facultad de Ciencias, UNAM} \\[0.6cm]

    \large \textbf{Tarea 5} \\[0.2cm]

    \begin{tabular}{l}
        \small Autor: Fernando Naed Ortega Montero\\
        \small Fecha: \today
    \end{tabular}
\end{center}
\[\] %borrar (Hypersnips meta.math declaration)


\section*{Preguntas conceptuales}

\begin{enumerate}[label=\alph*)]
\item Una partícula se encuentra inicialmente en la región $ a<x<b $. Si después de cierto tiempo la probabilidad de encontrarla en esa región disminuye, ¿significa esto que se ha perdido probabilidad? 
\item Suponga que al medir una observable $ \hat{M} $ sobre un sistema se obtiene el eigenvalor $ m_{i}$. ¿Qué sucede con el estado del sistema después de la medición y con qué probabilidad se obtendrá, al medir inmediatamente de nuevo la misma observable $ \hat{M} $, un eigenvalor $ m_{j} $ diferente de $ m_{i} $.
\item ¿Cuáles son las consecuencias de que la ecuación de Schr\"odinger sea una ecuación diferencial de segundo orden en la posición y de primer orden en el tiempo?
\item Un sistema se encuentra inicialmente en un estado $ \ket{\psi} $. Si el sistema no es sometido a ninguna medición, ¿puede cambiar su estado con el tiempo?.

\end{enumerate}

\section*{Problemas}


\begin{enumerate}
\item El momento de los Hermitianos

\begin{enumerate}
\item Demuestre que el operador derivada en la representación de posiciones, $ \hat{D_{x}} \doteq \frac{\partial }{\partial x}  $, no es Hermitiano. (Calcule $ \braket{\psi|\hat{D_{x}}|\phi} $ usando las funciones de onda $ \psi(x) \braket{x|\psi} $ y $ \braket{x|\phi} $ cuadrado integrables.) 

\item Usando el resultado anterior, muestre que el operador de momento lineal representado en el espacio de posiciones como $ \hat{p_{x}} \doteq -i\hslash\hat{D_{x}} $ sí es Hermitiano en la representación de posiciones.

\end{enumerate}

\item ¿Dónde está la partícula?
En el instante $ t = 0 $ una partícula está representada por la función de onda,

\[
\psi(x,0) = \begin{cases}
A(x/a), & 0 \leq x \leq a, \\
A(b/x)/(b-a), & a \leq x \leq b, \\
0, & en otro caso
\end{cases}
\]

donde $ A $, $ a $ y $ b $ son constantes positivas.


\begin{enumerate}
\item Normalice $ \psi(x,0) $ (es decir, encuentre $ A $ en una función de $ a $ y $ b $).
\item Grafique $ \psi(x,0) $ como función de x.
\item ¿Dónde es má probable encontrar a la partícula en $ t=0 $?
\item ¿Cuál es la probabilidad de encontrar la partícula a la izquierda de a? Compruebe su resultado en los casos límite $ b=a $ y $ b=2a $.
\item Calcule la integral

\[
\langle \hat{x} \rangle = \int_{-\infty}^{\infty} \psi*(x, 0)x\psi(x, 0) \, dx
\]

La cantidad $ \langle \hat{x} \rangle  $ se denomina valor esperado del operador de posición $ \hat{x} $ y está relacionada con el calor promedio de la posición que se obtiene al realizar muchas mediciones sobre sistema preparados de la misma manera.


\end{enumerate}

\item A la rapidez de la luz y hasta ahí.
La energía de una partícula relativista está dada por,

\[
E^{2} = m^{2} c^{4} + p^{2} c^{2},
\]

donde $ m $ es la masa de la partícula, $ p $ la magnitud del momento lineal y $ c $ la rapidez de la luz en el vacío. Obtenga, mediante un procedimiento análogo al visto en clase, la "versión relativista" de la ecuación de Schr\"odinger, denominada ecuación de Klein-Gordon. Para ello, parta de que $ \psi(x,t) $ es una onda plana y utilice la expresión relativista de E y las relaciones de Planck y de De Broglie. 


\end{enumerate}


\section*{Problemas}




\end{document}